% SPDX-License-Identifier: Apache-2.0
% Technical graph documentation; included by the main compendium.

\section{Stabile Ziele, kontextabhängige Struktur}
\label{sec:context-model}

Der Graph beantwortet, \emph{was} eine lernende Person können soll und
welche fachlichen Voraussetzungen dafür benötigt werden. Ein Bildungsgang
beantwortet zusätzlich, \emph{wo} diese Ziele verortet sind. Ein persönlicher
Lernplan bestimmt, \emph{wann} daran gearbeitet werden soll. Diese drei
Fragen brauchen unterschiedliche Daten: Das Umordnen eines Lehrplans darf
weder eine neue Kompetenz erzeugen noch bereits gespeicherte Ergebnisse
löschen.

Das allgemeine Zielmodell unterscheidet vier semantische Ebenen
und eine Darstellungsebene (Tabelle~\ref{tab:layers}). Es beschreibt die Richtung der schrittweisen
Migration; nicht jede hier genannte Datenstruktur ist damit bereits in
allen Landschaften vollständig umgesetzt.

\begin{table}[htbp]
\centering\small
\begin{tabularx}{\linewidth}{@{}P{27mm}XX@{}}
\toprule
Ebene & Bedeutung & Beispiel und Grenze \\
\midrule
Ziele & Konkrete Kompetenzen und fachliche Bündel; hier liegen
\code{requires} und \code{contains}. & Eine lineare Gleichung lösen;
Grundlagen linearer Funktionen. \\
Programme & Jahre, Semester, Module, Phasen und andere strukturelle
Einheiten. & Jahrgang 8 ist ein organisatorischer Abschnitt, keine
Kompetenz. \\
Platzierungen & Kontextuelle Zuordnung eines Ziels zu einer
Programmeinheit. & Dasselbe Ziel wird in einem Bildungsgang in Jahrgang 8,
in einem anderen in Jahrgang 9 erworben. \\
Kompetenzachsen & Fachliche Taxonomien und querschnittliche
Kategorien. & K1 \glqq Mathematisch argumentieren\grqq{} ist eine Kategorie;
\glqq Argumentationslücken erkennen\grqq{} kann ein konkretes Ziel sein. \\
Kompositionssicht & Geprüfter Aufbau des sichtbaren Baums für einen
aufgelösten Bildungsumfang. & Eine Landesansicht verweist auf kanonische
Teilbäume, ohne deren Ziele zu kopieren. \\
\bottomrule
\end{tabularx}
\caption{Fünf Modellierungsaufgaben, ein gemeinsamer Bestand stabiler Ziele.}
\label{tab:layers}
\end{table}

\subsection{Identität trotz unterschiedlicher Einordnung}

Die Identitäts-, Attribut- und Atomaritätsregeln aus
Abschnitt~\ref{sec:core} gelten in allen Schichten unverändert.
Eine neue Platzierung erzeugt deshalb keine neue Kompetenz.
Ein Programmabschnitt vom Typ \code{exam} ist noch kein
Prüfungsziel mit Aufgaben, Bewertungsraster und Lernergebnis.
Und \code{phase: Q1} erzeugt weder eine Elternkante zu einem
Q1-Knoten noch eine Voraussetzung zwischen Q1 und Q2.

Historische Bäume können Jahrgangs- und Phasenknoten als
Kompatibilitätsgerüst enthalten. Im allgemeinen Zielmodell
übernimmt explizite Programmstruktur diese Aufgabe; die
fachlichen Ziele bleiben eigenständig. Der gleiche Inhalt
kann so in mehreren Bildungsgängen anders eingeordnet werden,
ohne seine Ziel-ID oder gespeicherte Ergebnisse zu kopieren.

\subsection{Platzierungen und Projektionen}

\code{programUnits} beschreiben im Zielmodell einen expliziten Baum oder
Wald von Programmeinheiten. \code{goalPlacements} verbinden Ziele mit
diesen Einheiten. Eine \code{primary}-Platzierung liefert den normalen
Programmanchor; \code{secondary} zeigt zusätzliche Relevanz,
\code{assessed} einen Prüfungsbezug. Sekundäre und prüfungsbezogene
Zuordnungen erzeugen keine weiteren sichtbaren Kopien desselben Ziels.

Für einen aufgelösten Kontext gilt: Genau eine passende primäre
Platzierung ergibt einen eindeutigen Anker. Ohne passende Platzierung
darf die Darstellung keinen Elternknoten erfinden. Mehrere passende
primäre Platzierungen benötigen eine ausdrücklich geprüfte Auflösung;
sonst ist die Projektion uneindeutig. Kontextmerkmale wie Bundesland,
Schulform, Stufe, Dauer und Kursprofil ändern dabei nicht die Ziel-ID.

Die normale Inhaltssicht folgt ausschließlich den expliziten
\code{contains}-Kanten. Die Programmsicht folgt Programmeinheiten und
Platzierungen. Die Kompetenzsicht folgt Katalogeinträgen und
Kompetenzreferenzen. Die Reihenfolge der Kinder kommt aus der dafür
deklarierten Datenreihenfolge; eine topologische Sortierung nach
Voraussetzungen ist eine gesondert benannte didaktische Ansicht.

\subsection{Geprüfte Kompositionssichten}

Eine Kompositionssicht beschreibt den sichtbaren oberen Aufbau eines
Curriculums durch Strukturknoten, Beschriftungen, Reihenfolge und
Referenzen auf kanonische Teilbäume. Sie enthält keine neu verfassten
Atomziele und keine neuen \code{requires}-Kanten. Wo eine solche Sicht
für den Kontext maßgeblich ist, darf ein Renderer sie nicht durch
Vermutungen aus Titeln, \code{phase} oder alternativen Platzierungen
überschreiben.

Der aufgelöste Standardbaum muss jedes sichtbare Ziel höchstens einmal
und unter höchstens einem sichtbaren Elternknoten zeigen. Ein globaler
DAG kann mehrere fachliche Elternbezüge besitzen; die sichtbare
Baumprojektion benötigt dafür eine deterministische, geprüfte
Auflösungsregel. Überlappende Teilbaumreferenzen sind ohne einen solchen
Vertrag ungültig. Eine Rollenregel für überlappende Referenzen ist noch
keine Erlaubnis, Ziele im sichtbaren Baum zu duplizieren.

Die Kompilation soll früh erfolgen: Kontext auflösen, passende Sicht
auswählen, kanonische Referenzen expandieren, Rollen und sichtbare Eltern
prüfen, dann das Ergebnis für die Laufzeit bereitstellen. Damit wird
Strukturvalidierung möglich, bevor eine Benutzeroberfläche daraus einen
Baum zeichnet.

Ein autoritativ als \code{curricularArea} klassifizierter Cluster
darf nicht über eine direkte \code{goalEntry}-Referenz als
undurchsichtiges, atomar wirkendes Laufzeitziel erscheinen. Eine
\code{canonicalSubtree}-Referenz erhält seine geprüften Nachkommen;
alternativ benennt eine explizite Sichtstruktur die beabsichtigten
Atomziele. Diese Darstellungsregel verändert keine Projektionsrolle.

\paragraph{Quellen.}
\srcref{docs/concept/skill-graph/general-goal-system-and-migration.md}{Allgemeines Zielsystem und Migration};
\srcref{docs/concept/skill-graph/view-projection-and-goal-placement.md}{Sichten und Zielplatzierung}.

\section{Persönlicher Umfang, Fokus und sichtbare Ziele}
\label{sec:personal-scope}

\subsection{Vier Personalisierungsstufen}

Die Auswahl \glqq Schule\grqq{}, \glqq Hochschule\grqq{} oder \glqq Sprachen\grqq{} ist zunächst
Katalognavigation. Erst danach beginnen die vier
Personalisierungsstufen. Sie sind von den curricularen Levels zu
unterscheiden, die Erwerbsetappen mit Voraussetzungen und Zeitbudgets
bezeichnen.

\begin{enumerate}
\item \textbf{Basiscurriculum:} ein ausdrücklich auswählbares
Wurzelcurriculum stellt den Zielbestand bereit.
\item \textbf{Personal Curriculum:} der längerfristig festgelegte,
backendseitig verantwortete Bildungsumfang. Dazu gehören beim
Schulcurriculum der Landes- beziehungsweise kanonische Kontext,
gegebenenfalls G8/G9, eine ausdrücklich bestimmte Stufe, die Fächer
und deren Kursprofile.
\item \textbf{Lernfokus und aktives Ziel:} Ein oder mehrere Wurzeln
begrenzen vorübergehend den Arbeitsbereich; genau ein atomarer
Lerngegenstand ist das aktuelle Ziel. Fokus und aktives Ziel sind
veränderlich.
\item \textbf{Mastery:} Das binäre Ergebnis \glqq geschafft\grqq{} oder
\glqq noch nicht geschafft\grqq{} gehört zur stabilen atomaren Ziel-ID.
Die Ergebnisse bleiben bei Änderungen der Sicht oder des
Fokus erhalten. Clusterfortschritt wird daraus abgeleitet.
\end{enumerate}

Ein Kursprofil wie LK ist ein Fachattribut und bestimmt nicht implizit
die Schulstufe. Fehlende Angaben zur Stufe bleiben ungeklärt. Eine
Änderung des Personal Curriculum prüft Fokus und aktives Ziel erneut
gegen die neue Zielmenge; sie darf globale Ergebnisse weder löschen
noch neu vergeben. Die aktuelle Konfiguration des Personal Curriculum
erfolgt über die dafür vorgesehenen SkillPilot-Oberflächen. Der
OpenAI-V1-Coach ändert diese Konfiguration nicht selbständig.

Alle \code{target}-Ziele im Personal Curriculum bleiben beabsichtigte
Lernziele. Ein enger Fokus reduziert den vorübergehenden Arbeitsbereich,
nicht den längerfristigen Gesamtumfang oder dessen Abschlussbedingungen.

\subsection{Zielrolle und Voraussetzungshilfe}

Eine Referenz in einer Kompositionssicht kann ausdrücklich eine von
zwei Rollen tragen:

\begin{description}
\item[\code{target}] Das Ziel gehört zur auswählbaren Lernmenge,
zur Frontier, zum Fortschritt und zum Abschluss dieser Sicht.
\item[\code{prerequisiteOnly}] Das Ziel bleibt für
Voraussetzungsprüfungen mit seiner globalen Identität und seinem
gespeicherten Zustand verfügbar. Es ist in dieser Sicht kein
auswählbares Ziel und zählt nicht zum Fortschrittsnenner.
\end{description}

Fehlt eine Rollenangabe, gilt \code{target}. Teilbaumreferenzen vererben
ihre Rolle an Nachkommen. Bei überlappenden Rollenfestlegungen entscheidet
die Spezifität: Eine direkte \code{goalEntry}-Referenz geht einer
Teilbaumvererbung vor; eine nähere Teilbaumwurzel geht einer breiteren
vor; erst bei gleicher Spezifität gewinnt \code{target}. Diese Regel
löst die Rolle, nicht eine ansonsten unzulässige doppelte Darstellung.

Rollen werden nicht aus Alter, Phase oder Voraussetzungspfaden erraten.
Ein mathematisches Ziel kann in Mathematik ein Lernziel und in Physik
eine unterstützende Voraussetzung sein. Eine gemeinsame Planung
identifiziert es dennoch nur einmal. Seine Physikrolle löscht keine
Mathematikergebnisse und erzeugt auch keinen zweiten Kompetenzerwerb.

Jede direkte kanonische Voraussetzung eines sichtbaren Ziels muss
in der maßgeblichen Level-2-Komposition als \code{target} oder
\code{prerequisiteOnly} auflösbar sein. Fehlt sie, ist die
Konfiguration fehlerhaft und das Ziel bleibt gesperrt. Nur bei
ausgelassenen Anforderungen, die von übergangsweise vorhandenen
Legacy-Clustern geerbt werden, bleibt das besondere
Kompatibilitätsverhalten bestehen. Diese Gültigkeitsprüfung ist
von der späteren Level-3-Fokuswahl getrennt: Eine gültig
aufgelöste Voraussetzung außerhalb des Fokus blockiert den
Normalmodus weiterhin nicht.

\subsection{Anwendbarkeit ist keine Platzierung}

\code{applicability} ist kompilierte Sichtmetadateninformation, etwa
für einen Bundeslandfilter. Sie wird aus überprüfbaren Quellen,
Mappings, Provenienz und begründeten Ausnahmen gewonnen. Sie beantwortet,
in welchem Kontext ein Ziel sichtbar ist. Eine Platzierung beantwortet,
wo es in einer Programmstruktur liegt. Keine der beiden Angaben ist
ein Beweis für individuelle Beherrschung.

Bei landesspezifischen Ergänzungen kann eine ausdrücklich deklarierte
Vererbungsgrenze verhindern, dass eine breite Quellenzuordnung eines
Vorfahren als Quellenbeleg aller neuen Nachkommen ausgegeben wird.
Direkte Quellen und Zuordnungen im Ergänzungszweig bleiben maßgeblich.
Ebenso kann die Sichtbarkeit eines lokalen Übungs- oder Prüfungsendpunkts
unter einem dafür vorgesehenen Vertrag aus seinen Voraussetzungen
abgeleitet werden. Das belegt seine Anwendbarkeit, nicht automatisch
seine Lehrplanabdeckung.

\subsection{Fokuswechsel ohne impliziten Kompetenzerwerb}

Ein Fokus besteht aus der Vereinigung der gewählten Wurzeln und ihrer
sichtbaren Nachkommen, eingeschränkt auf das Personal Curriculum.
Im Normalmodus werden die effektiven Voraussetzungen innerhalb dieses
Fokus geprüft; der explizite strenge Modus prüft global. Eine außerhalb
liegende Voraussetzung wird im Normalmodus weder zum Hindernis noch
dadurch gemeistert. Die formalen Unterschiede der Frontier-Berechnung
werden im mathematischen Teil behandelt.

Eine Fokusausweitung folgt dem eindeutigen sichtbaren Elternpfad. Die
Backendoption muss eine echte Obermenge liefern und mindestens ein
noch offenes Ziel hinzufügen. Bei mehreren Wurzeln ersetzt sie nur den
abgedeckten Zweig und bewahrt die unabhängigen Wurzeln. Der Coach
übernimmt die veröffentlichte Option vollständig und wartet vor der
Änderung auf die Zustimmung der lernenden Person.

Eine automatische Ausweitung wird erst nach tatsächlichem Abschluss
des aktuellen Fokus vorgeschlagen. Eine leere Frontier kann auch einen
blockierten, unvollständigen Fokus bedeuten. Sie beweist keinen
Abschluss. Nach der Ausweitung werden neu enthaltene, unbeherrschte
Voraussetzungen regulär geprüft: Die Beherrschung eines abhängigen Ziels
überträgt keine Mastery rückwärts. Ihre Erforderlichkeit wird aus dem
vollständigen Bedarfskontext neu berechnet; rein implizit erledigte
Erweiterungen begründen keinen neuen Pflichtvorschlag.

\paragraph{Quellen.}
\srcref{docs/concept/levels-of-personalization.md}{Personalisierungsstufen};
\srcref{docs/concept/skill-graph/view-projection-and-goal-placement.md}{Projektionsrollen und Sichtvertrag}.

\section{Zieltypen und die Bedeutung eines Ergebnisses}
\label{sec:goal-types}

Der Graph erhält einen zieltypgerecht aufgelösten Zustand. Er entscheidet
nicht selbst, ob eine Antwort fachlich überzeugend war, ein
Orientierungsgespräch angemessen abgeschlossen wurde oder alle
Karten eines Decks sicher abgerufen wurden. Diese Entscheidungen
gehören zum jeweiligen Ergebnisvertrag.

Das Ergebnis jedes Lernziels bleibt binär: geschafft oder noch nicht
geschafft. Der Zieltyp bestimmt, wann es als geschafft gilt.
Bei einem Inhaltsziel ist dies die fachlich festgestellte Beherrschung;
bei Orientierung der abgeschlossene motivierende Einstieg, kein
fachliches Zertifikat. Bei Memorierung wird das eine Zielergebnis aus
den Kartendaten abgeleitet. Unterschiede der Abschlussbedingungen
erzeugen keine zusätzlichen Mastery-Stufen oder parallelen Zielergebnisse.

\subsection{Verstehen und Anwenden}

Ein gewöhnliches Inhaltsziel beschreibt konzeptuelles oder prozedurales
Verstehen. Die Beurteilung muss die in der Beschreibung verlangten
Aspekte abdecken. Eine gelöste Aufgabe ist ein möglicher Beleg, aber
nicht schon per Definition die Beherrschung des ganzen Ziels.
Mehrschrittiger Transfer in einer Aufgabe kann ausreichend sein; eine
pauschale Mindestzahl von Aufgaben wäre ebenso unbegründet wie das
automatische Anerkennen jeder richtigen Einzelantwort.

Im coachgeführten Ablauf wird eine ausreichend belegte Beherrschung
unmittelbar gespeichert. Erst eine bestätigte Speicherung erlaubt
die Aussage, das Ergebnis sei gesichert. Eine zusätzliche Zustimmung
zum Speichern ist kein fachlicher Nachweis und wird für den regulären
Erfolgsfall nicht als weitere Abschlussbedingung eingeführt.
Der anschließende Themenwechsel bleibt eine andere Entscheidung:
Der Coach bietet das backendseitig gewählte nächste Thema an und wartet
auf erkennbares Einverständnis, bevor neues Material beginnt.

\subsection{Orientierung als eigener Abschlussvertrag}

Ein autoritativ als \code{orientation} klassifiziertes Ziel zeigt
konkrete Möglichkeiten und ehrliche positive Perspektiven des
folgenden Lernstoffs. Es prüft weder Vorwissen noch Fachbegriffe,
Rechenfertigkeiten, Abruf, Transfer oder Prüfungsleistung. Wie bei allen
Lernzielen ist das Ergebnis binär; eine Prozentbewertung der persönlichen
Motivation wäre zudem inhaltlich irreführend.

Wählt die lernende Person eines der angebotenen Interessen, beginnt
dadurch zunächst eine passende motivierende Vertiefung. Der Coach
verknüpft das Interesse mit etwas Konkretem, das die Person verstehen,
erkunden, gestalten oder tun kann, und lädt zu einer einfachen
persönlichen Reaktion, Wahl oder Frage ein. Das bloße Wiederholen
eines Angebotslabels schließt die Orientierung nicht ab.

Nach einer sinnvollen Antwort auf diese Vertiefung oder dem
ausdrücklichen Wunsch, die Orientierung zu verlassen, folgt eine
positive Zusammenfassung mit einem gesonderten Abschlussangebot.
Erst die weitere Zustimmung dazu begründet den Abschlussmarker.
Dieser Ablauf darf weder als Fachprüfung umgedeutet noch auf normale
Inhaltsziele übertragen werden. Orientierungsknoten sind reine
voraussetzungsfreie atomare Einstiege; Prüfungs- und Memory-Metadaten
gehören nicht auf denselben Knoten.

Bei einem echten Neueinstieg ohne bisherigen Fortschritt im projizierten
Umfang sind offene Orientierungsziele die einzigen auswählbaren Frontier-Ziele.
Lernende mit bereits bestehendem Fachfortschritt werden nicht auf diese
Einstiegsregel zurückgesetzt.

\subsection{Memorierung: Übung und verifizierter Abruf}

Ein Memorierziel bündelt die erforderlichen Karten eines Decks.
Die Karten besitzen zwei getrennte Prozesszustände. Die SRS-Übung steuert
Wiederholungsintervalle. Der verifizierte Abruf hält fest, ob eine
Karte in einer kontrollierten Abfrage korrekt beantwortet wurde.
Das Umdrehen einer Karte oder eine Selbsteinschätzung in der Übung
ist kein solcher Nachweis. Aus den Kartendaten folgt genau ein binäres
Mastery-Ergebnis des Lernziels; Übungs- und Abrufprozess sind keine
zusätzlichen gespeicherten Zielergebnisse.

Sei $K_g$ die aufgelöste Menge erforderlicher Karten eines Ziels $g$.
Für einen Zeitpunkt $t$ bezeichnet $\Due(k,t)$ die Fälligkeit und
$\Verified(k)$ einen aktuell gültigen bestandenen Abruf. Der
\emph{dokumentierte Zieltypvertrag} verlangt beide Bedingungen:
\[
  \Complete_{\mathrm{memory,doc}}(g,t)
  \iff
  \bigl(\forall k\in K_g:\neg \Due(k,t)\bigr)
  \land
  \bigl(\forall k\in K_g:\Verified(k)\bigr).
\]
Diese Formel setzt ein vollständig aufgelöstes, gültiges Deck voraus.
Fehlende Kartendaten werden dadurch nicht als leeres, automatisch
abgeschlossenes Deck interpretiert.

\textbf{Abweichung des geprüften Codes.} Der zum Dokumentstand
gelesene \code{LearnerService} und der Frontend-Hook
\code{useSrsMastery} leiten Memorier-Mastery dagegen allein aus dem
verifizierten Abruf aller erforderlichen Karten ab. Für ein nichtleeres,
vollständig aufgelöstes Deck lautet diese Implementierungsregel:
\[
  \Complete_{\mathrm{memory,code}}(g)
  \iff \forall k\in K_g:\Verified(k).
\]
Neue SRS-Fälligkeit entzieht danach nicht automatisch die Mastery.
Ein fehlgeschlagener erneuter Abruf kann sie hingegen zurücknehmen.
Dokumentierte Sollsemantik und aktueller Code sind damit nicht
deckungsgleich; dieses Kompendium trifft keine stillschweigende
Migrationsentscheidung. Die Abweichung ist einmal im Zieltypvertrag
aufzulösen. Die allgemeine Abschlussberechnung verwendet dasselbe binäre
Zielergebnis wie die übrigen Oberflächen; es führt keine eigene
Anerkennung von Memorier-Mastery ein. Fällige Übung bleibt eine
Prozessaussage über die Karten.

Der gegenwärtig dokumentierte SRS-Vertrag verwendet SM-2-Zustand mit
Intervall, Wiederholungszähler, Leichtigkeitsfaktor und nächstem
Wiederholungsdatum. Eine Karte ohne gespeicherten Zustand ist fällig.
Erfolgreiche Übung verlängert nach den festgelegten Regeln das
Intervall; unzureichende Übung setzt die Wiederholungsfolge zurück.
Diese Zeitplanung ersetzt den verifizierten Abruf nicht.

Eine harte Abfrage kann die lernende Person beim aktiven Memorierziel
auch beginnen, wenn die SRS-Warteschlange noch nicht leer ist. Zuerst
wird nur die Frage abgerufen, dann antwortet die Person ohne
Kartenhilfe. Erst anschließend wird die erwartete Antwort geladen
und der Abruf als bestanden oder fehlgeschlagen gespeichert.
Ein bestandener Abruf aktualisiert auch die SRS-Spur als starke
erfolgreiche Wiederholung. Ein Fehlschlag hebt den verifizierten
Status dieser Karte auf und führt sie in die Übung zurück.

\subsection{Prüfungsziele im geltenden Ablauf}

Prüfungsziele verwenden \code{examData} mit Aufgabeninhalt,
Lösungsinhalt und Bewertungsraster. Das Raster enthält maximale
Punkte, Bestehensgrenze und bewertbare Schritte. Für gültig
vergebene Schrittpunkte $p_j$ gilt im aktuellen Vertrag:
\[
  \Score=\min\!\left\{\sum_j p_j,\;\MaxPoints\right\},
  \qquad
  \Passed\iff \Score\ge \PassingPoints.
\]
Die Punktevergabe folgt den einzelnen Rastervorgaben. Die
Gesamtsummenformel allein legitimiert keine frei erfundenen
Punkte, negativen Bewertungen oder Änderungen des Aufgabeninhalts.

Der normale Lernzugang prüft direkte effektive Voraussetzungen einschließlich
der atomaren Erfüllung erforderlicher Cluster. Für Prüfungen gilt zusätzlich
der allgemeine direkte Zugang: Fehlende Vorbereitungsmastery sperrt einen
bewusst gewählten Prüfungsversuch nicht. Aufgabenverfügbarkeit, persönlicher
Geltungsbereich und autorisierte Sitzung bleiben erforderlich. Dieser Zugang
ist lokal implementiert; tatsächliches Deployment und Coach-Abnahme sind
getrennt nachzuweisen.

Im Prüfungsmodus erhält die Person die vollständige, unveränderte Aufgabe
und bearbeitet sie ohne inhaltliche Hilfen. Ihre vollständige Einreichung
wird anhand des Rasters bewertet. Bei Erfolg wird Mastery unmittelbar
gespeichert. Ein Fehlversuch speichert weder Ergebnis noch Punkte und lässt
vorhandene Mastery bestehen. Erst die bestätigte Speicherung erlaubt die
Aussage, der Erfolg sei gesichert.

Nach der Bewertung erhält die Person Ergebnis, Punkte, Begründung und
Lösung; neues Material beginnt nach ihrer weiteren Entscheidung. Die
unbegrenzte Wiederholung derselben Aufgabe bleibt möglich. Bekannte Lösungen
können die Aussagekraft eines Versuchs verringern, erzeugen aber keine zweite
Anerkennungsstufe: Vorhandene Prüfungs-Mastery zählt unmittelbar in der
allgemeinen Abschlussberechnung. Entfallende Arbeit wird aus dem vollständigen
Bedarfsgraphen und der Mastery aller Zielarten abgeleitet, nicht ausschließlich
aus bestandenen Prüfungen.

\paragraph{Quellen.}
\srcref{docs/concept/skill-graph/node-types.md}{Zieltypen und Prüfungsablauf};
\srcref{AGENTS.md}{Projektregeln zu Orientierung, Evidenz und Fortsetzung};
\srcref{backend/src/main/java/com/skillpilot/backend/service/LearnerService.java}{Backend-Ableitung von Memorier-Mastery};
\srcref{app/src/hooks/useSrsMastery.ts}{Frontend-Ableitung von Memorier-Mastery}.

\section{Zeitplanung als ergänzende Sicht}
\label{sec:time-context}

Ein Zeitplan bindet einen erklärten Zielumfang, eine Curriculumfassung,
Kalender, verfügbare Kapazität und Planungsversion. Absolute Termine
gehören nicht auf kanonische Ziele, Cluster oder Platzierungen.
Derselbe Inhalt kann in verschiedenen Kursen zu verschiedenen
Zeiten bearbeitet werden, ohne seine fachliche Identität zu ändern.

Die Zeitplanung ist eine zusätzliche Ausführungsplanung und keine
fünfte Personalisierungsstufe. Sie darf Ziele priorisieren, die im
aktuellen Arbeitsbereich verfügbar sind. Sie erzeugt keine
Voraussetzungskanten, keine neue Frontier und keine Lernergebnisse.
Ein curricularer Level ordnet Erwerbsetappen gröber; ein datierter
Ablaufplan muss darüber hinaus gewährleisten, dass eine Voraussetzung
rechtzeitig vor ihrer geplanten Verwendung erarbeitet wird.

Verfügbare Kapazität berücksichtigt tatsächliche Unterrichts- oder
Lernminuten nach Ferien, Ausfällen und anderen Ausschlüssen.
Aktivitäten, geschützte Reserve und verbindliche Fachbudgets
müssen darin ohne Doppelzählung Platz finden. Die formale
Budgetprüfung und ihre Anforderungen an kalibrierte
Aufwandsschätzungen folgen in Abschnitt~\ref{sec:levels}.

Plan-Soll und tatsächliche Zielerreichung verwenden denselben
benannten Umfang. Gespeicherte Lernergebnisse sind die Quelle
der Zielerreichung; eine zusätzliche manuelle
Unterrichtsabdeckung begründet oder verbirgt keinen Abschluss.
Fehlende Ergebnisse bleiben als nicht verfügbar erkennbar.
Neue Planrevisionen dürfen historische Ziele, Ergebnisse und
frühere Aussagen nicht stillschweigend umdeuten.

Der dokumentierte lokale Planungspilot und die angestrebte
backendseitig maßgebliche Architektur sind getrennt zu betrachten.
Insbesondere behauptet die Leveldefinition noch keine
implementierte fachübergreifende Budgetprüfung.

\paragraph{Quelle.}
\srcref{docs/concept/didactic/curriculum-time-axis-and-pacing.md}{Zeitachse, Kapazität und Planungsgrenzen}.

\section{Herkunft, Materialien und belastbare Qualitätsaussagen}
\label{sec:evidence-context}

\subsection{Drei Arten von Verweisen}

Curriculare Herkunft, eigene Veranschaulichung und optionales
Lernmaterial haben verschiedene Funktionen. Ein offizieller
Quellenverweis begründet die Aufnahme oder Formulierung eines
Lernziels. Ein eigenes Lernzielbild veranschaulicht dessen Kernidee.
Ein Buch, Video oder eine Simulation kann bei der Erarbeitung
helfen. Keiner dieser Verweise ist für sich ein individueller
Kompetenznachweis.

Die verbindliche Inhaltsarchitektur lautet: Materialien
referenzieren das Curriculum; das Curriculum wählt keine
Materialanbieter. Neue optionale externe Lernmaterialien gehören
in getrennte Inhaltspakete, die auf stabile Ziel-IDs verweisen.
Die Auswahl eines Pakets verändert weder Personal Curriculum
noch Voraussetzungen, Mastery oder Abschlussregeln. Unvollständige
Materialabdeckung darf normale Lernarbeit nicht sperren.

Offizielle Quellen, Mappings und strukturierte Provenienz bleiben
am Curriculum nachvollziehbar. Eigene Zielvisualisierungen behalten
ihren gesonderten kanonischen Referenz- und Qualitätsvertrag.
Bestehende externe Links auf kanonischen Zielen sind gegebenenfalls
Migrationsbestand; daraus folgt keine dauerhafte zweite
Autorenarchitektur. Eine reine Linklieferung belegt außerdem
nicht, dass ein Coach den verknüpften Inhalt gelesen hat.

\subsection{Quellenbegründung für Ziele und Beziehungen}

Eine überprüfbare Herkunftskette verbindet das kanonische Ziel mit
einer Quelle, einem konkreten Quellabschnitt, einer nachvollziehbaren
Extraktion und einer geprüften Mappingentscheidung. Diese kann
eins zu eins, teilweise, aggregierend oder auf mehrere Ziele
aufteilend sein. Eine verständliche Begründung erklärt diese Form,
anstatt jede Zuordnung als wortgleiche Lehrplanübernahme darzustellen.

Für \code{requires} und \code{contains} ist ebenfalls zu
unterscheiden, ob eine Beziehung ausdrücklich in der Quelle steht
oder eine fachlich begründete Modellierungsentscheidung ist.
Ein plausibler Zieltext erlaubt keine nachträglich erfundene
Provenienz. Ungeklärte Quellen, veraltete Fassungen und
widersprüchliche Zuordnungen bleiben sichtbare Befunde.

Die menschenlesbare Begründung ist eine Darstellung vorhandener
Evidenz, kein zweites führendes Register. Strukturierte Referenzen
und Fingerprints ermöglichen, dieselbe Begründung in Dokumenten,
Exporten und Laufzeitwerkzeugen konsistent auszugeben. Ein
zusätzlicher MEM/FWU-Abgleich kann die klassische Quellenroute
stützen oder infrage stellen; er ersetzt sie nur unter einem
ausdrücklich geprüften Vertrag für den konkreten Umfang.

\subsection{Algorithmische Gültigkeit und fachliche Qualität}

Die Ebenen der Qualität bauen aufeinander auf, sagen aber nicht
dasselbe aus. Ein Parser erkennt fehlerhaftes JSON; eine
Referenzprüfung unbekannte IDs; ein Zyklustest unerlaubte
Kantenfolgen. Keiner dieser Tests beurteilt allein, ob die
Kompetenz sinnvoll geschnitten, eine Erklärung verständlich
oder eine Aufgabe didaktisch geeignet ist.

Die in Abschnitt~\ref{sec:validity} erläuterten
Gültigkeits- und Routenprofile brauchen jeweils einen benannten
Prüfumfang. Ein lokales Migrationsgate oder ein einzelner
grüner Strukturcheck bestätigt nicht automatisch sämtliche
Anforderungen an Graph, Projektion und Lernbarkeit.

Für vertiefte Curriculum-QS unterscheidet SkillPilot fünf
gleichzeitig aktuelle Nachweisbereiche:
\begin{description}
\item[D -- Beschreibung:] Ist der Zieltext fachlich korrekt und klar?
\item[P -- positive Verständnisnachweise:] Welche beobachtbare
Leistung zeigt das geforderte Verständnis?
\item[A -- Atomarität:] Ist die Kompetenz sinnvoll zugeschnitten?
\item[M -- Memory:] Ist die Entscheidung über notwendige
Memorieranteile begründet und aktuell?
\item[V -- Visualisierung:] Ist die Darstellung fachlich geeignet
und an die geprüften Bild- und Zielstände gebunden?
\end{description}

Der Reifegrad M7 bezeichnet den vollständigen geltenden
maschinellen Abschluss dieser Anforderungen und der
vorausgesetzten unteren Reifegrade für den erklärten Umfang.
Er ist keine menschliche Freigabe, keine menschliche Erprobung
und kein Nachweis von Lernwirksamkeit. Solche Aussagen
benötigen eigene Prüfungen und eigene Evidenz. Auch mehrere
übereinstimmende KI-Reviews ersetzen keine gesondert
vorgeschriebene menschliche Entscheidung.

Eine Prüfung bindet die genaue Fassung ihrer Gegenstände.
Ein gültiges Bildreview für einen alten Asset-Hash bestätigt
nicht das neue Bild. Eine Änderung von Zieltext oder
Voraussetzungen kann bisherige inhaltliche Evidenz berühren;
ein rein technischer Umzug kann unter einem eng begründeten
Kompatibilitätsnachweis anders zu behandeln sein. Beides darf
nicht durch pauschales Erneuern von Hashes verborgen werden.

\subsection{Veröffentlichung und Austausch}

Fachliche Änderungen erfolgen in kanonischen Autorenquellen.
Pakete, Bücher, Bilder und Validierungsberichte werden daraus neu gebaut;
historische veröffentlichte Archive bleiben unverändert.

JSON- und Ontologiepakete sollen aus demselben fachlichen Stand entstehen.
Das JSON-Laufzeitmodell ist die normative Referenz. Der Ontologie-Roundtrip
muss Zielidentitäten, Relationen, Rollen, Texte und weitere fachliche Felder
gemäß ihrem Gleichheitsvertrag erhalten. Unterschiedliche JSON-/RDF-Texte
benötigen semantische Gleichheit, binäre Ressourcen identische Bytes.
Das Austauschformat darf keine abweichende Laufzeitsemantik einführen.

Externe Zielreferenzen müssen vollständig aufgelöst sein. Ein Physikpaket
mit mathematischen Voraussetzungen benötigt deren Definitionen;
leere Platzhalter genügen nicht. Diese vollständige fachliche Closure
bestimmt noch nicht, welche Ziele für die Person als \code{target}
auswählbar sind.

Veröffentlichungs- und Prüfaussagen nennen Artefakt, Umfang, Version und
gebundene Quellen. Ein lokaler Build belegt weder fremde Hostanzeige noch
pädagogische Wirksamkeit. Dieses Kompendium erteilt keine Release- oder
Produktionsfreigabe.

\subsection{Lizenzen und private Lernergebnisse}

Software, technische Schemas, Entwicklerdokumentation und funktionale
Coachanweisungen stehen unter Apache-2.0; eigene Skill-Landschaften,
Zieltexte, didaktische Aufgaben, Karten, Bilder und Bildungsdokumentation
unter CC-BY-4.0. Entscheidend ist die Funktion, nicht die Dateiendung.
Gemischte Sammlungen sind nicht pauschal doppelt lizenziert.

Fremde Quellen, Abhängigkeiten, Medien und Marken behalten ihre Bedingungen.
Quellenangabe oder KI-Herkunft sind weder Rechtefreigabe noch Qualitätsbeweis.
Exporte führen zutreffende Lizenzen, Herkunftsangaben und Änderungshinweise
mit. Private Lernergebnisse, Sitzungsdaten und Gespräche gehören weder zur
öffentlichen Inhaltslizenz noch in öffentliche Quellen- oder Reviewpakete.

\paragraph{Quellen.}
\srcref{docs/concept/skill-graph/source-and-resource-links.md}{Quellen und Materialien};
\srcref{docs/concept/skill-graph/human-readable-source-rationales.md}{Menschenlesbare Quellenbegründungen};
\srcref{docs/qa-ci/graph-validation-rules.md}{Algorithmische Graphvalidierung};
\srcref{docs/qa-ci/curriculum-quality-maturity-and-routes.md}{Reifegrade und Qualitätsregeln};
\srcref{docs/concept/skill-graph/learning-goal-book-and-evidence-review-pipeline.md}{Lernzielbuch und Evidenzreview};
\srcref{docs/concept/skill-graph/dual-curriculum-package-releases.md}{Curriculum-Pakete};
\srcref{LICENSING.md}{Lizenzaufteilung}.
